\scanpage{515}

We will have, in terms of $\alpha,\beta,f$, the following expressions for the action of $U$ on $\mathcal{G}$:
\[
\left\{\begin{aligned}
U(\rho)&=\operatorname{int}(f^{-1}\alpha)\deleted{\unsure{\operatorname{int}(\alpha)}}\rho
 =\operatorname{int}(f^{-1})(\zeta\rho)\\
U(\sigma)&=\operatorname{int}(f^{-1}\beta)\deleted{\unsure{\operatorname{int}(\beta)}}(\sigma)
 =\operatorname{int}(f^{-1})(\eta\sigma).
\end{aligned}\right.\tag{59}
\]
\deleted{\unsure{and}\ $\unsure{\zeta|\xi},\unsure{\eta}$\unsure{.}}

\uline{\unsure{Remark|Remark}} Finally, in all these developments, we have not really used the so-called ``fundamental hypothesis'' $(20)\deleted{\unsure{b}}$, \deleted{\unsure{and}} except part a). In the absence of b), one should introduce
\[
\begin{gathered}
\Sigma_{\rho,\sigma}^{\mathrm{eff}}\subset\Sigma_{\rho,\sigma}\subset\mathcal{G}\times\mathcal{G}\\[-1pt]
\begin{array}{c@{\qquad}l}
\rotatebox{90}{$=$}&\\[-3pt]
\multicolumn{2}{l}{\left\{(\zeta,\eta)\in\unsure{\mathcal{G}|\Sigma_\rho}\times\mathcal{G}\mid
\exists u\in B_{\rho,\sigma},\ u(\rho)=\zeta\rho,\ u(\sigma)=\eta\sigma\right\},}
\end{array}
\end{gathered}\tag{60}
\]
which is thus nothing other than the image of the map $(9)$: \uline{effective} pairs $(\zeta,\eta)$ satisfying the conditions defining \deleted{\unsure{the}\ $\unsure{(5)}$}\inserted{$\Sigma_{\rho,\sigma}$ (cf. $(8)$)}\unsure{)}

\deleted{The map
$\diagram{\begin{tikzcd}[ampersand replacement=\&,column sep=large]
S\unsure{\mathscr{G}|\mathcal{G}}_{\rho,\sigma}\arrow[r]\&\Sigma_{\rho,\sigma}
\end{tikzcd}}$
\inserted{$(28)$} (will not necessarily be surjective, its image \unsure{is|will be} $\Sigma_{\rho,\sigma}^{\mathrm{eff}}\inserted{\simeq B_{\rho,\sigma}}$), otherwise there is nothing to change.}
\margin{\unsure{then|by}\par\unsure{bijection|bis}}

let us write
\[
\begin{aligned}
S\unsure{\mathscr{G}|\mathcal{G}}_{\rho,\sigma}^{\mathrm{eff}}
 &=\text{Inverse image of }\Sigma_{\rho,\sigma}^{\mathrm{eff}}\\
 &\text{under }\diagram{\begin{tikzcd}[ampersand replacement=\&,column sep=large]
S\unsure{\mathscr{G}|\mathcal{G}}_{\rho,\sigma}\arrow[r,no head]\&\Sigma_{\rho,\sigma}
\end{tikzcd}}\quad(28),
\end{aligned}\tag{61}
\]
\deleted{the inverse image}
and it is on $S\unsure{\mathscr{G}|\mathcal{G}}_{\rho,\sigma}^{\mathrm{eff}}$ that we will have a structure

\mspage{514}

\scanpage{516}
\mspage{561}

of a group \unsure{isomorphic} \unsure{to|:} $S_0\mathcal{T}_{\rho,\sigma}^{0}$, by the isomorphism (which replaces (38))
\diagram{\begin{tikzcd}[ampersand replacement=\&,column sep=large]
 S_0\deleted{\unsure{\mathcal{T}}}\Gamma_{\rho,\sigma}^{0} \arrow[r,phantom,"\simeq"] \& S_0\unsure{\mathscr{G}|\mathcal{G}}_{\rho,\sigma}^{\mathrm{eff}}\,. \\
 (\underline{u},\underline{a},\underline{b}) \arrow[r] \& (\underbrace{ua^{-1}}_{\zeta},\underbrace{ub^{-1}}_{\eta})
\end{tikzcd}\quad(62)}

This being so, let us return to $M_{\rho,\sigma}^{\sim}$, and the exact sequence
\diagram{\begin{tikzcd}[ampersand replacement=\&,column sep=small]
 1 \arrow[r] \& \mathcal{G}_{0,3}^{+\wedge} \arrow[r] \& M_{\unsure{0,3|\rho,\sigma}}^{\sim} \arrow[r] \& \underbrace{\Gamma_{0,3}^{\sim}}_{\Gamma_{\mathbb Q}^{\sim}} \arrow[r] \& 1
\end{tikzcd}\quad(63).}

Consider the subgroup
\[
\begin{aligned}
(64)\qquad
\Gamma_{\mathbb Q}^{\prime\sim}
 &= \operatorname{Ker}\bigl(\Gamma_{0,3}^{\sim}
 \xrightarrow{(\varepsilon_3,\deleted{\unsure{\varepsilon_2}})}
 \{\pm1\}\deleted{\unsure{\times\{\pm1\}}}\bigr)\\
 &= \{\gamma\in\Gamma_{0,3}^{\sim}\mid\mu(\gamma)\equiv1\ (\unsure{3}\deleted{\unsure{12}})\}.
\end{aligned}
\]
\deleted{\illegible{3}}
\deleted{\unsure{two factors}}

and let \deleted{$\unsure{M_{0,3}^{\prime\sim}}$}
\[
 (65)\qquad M_{0,3}^{\prime\sim}
 = \text{\unsure{Inverse image} of }\Gamma_{0,3}^{\prime\sim}\text{ in }M_{0,3}^{\sim}.
\]

Let $\mathcal{G}$ be a profinite group, and
\[
 (66)\qquad\psi:\mathcal{G}_{0,3}^{+\wedge}\longrightarrow\mathcal{G},
\]
a \inserted{continuous} surjective homomorphism --- which amounts to the same thing as giving two elements
\[
 (67)\qquad\rho_{\mathcal{G}},\unsure{\sigma_{\mathcal{G}}}\in\mathcal{G}
\]
(the images of $\rho,\sigma$ under $\psi$) satisfying
\[
 (68)\qquad\rho_{\mathcal{G}}^3=\sigma_{\mathcal{G}}^2=1.
\]
\mspage{515}

\scanpage{517}

\margin{introduce it by equivalent conditions:

\unsure{a)}
\diagram{\begin{tikzcd}[ampersand replacement=\&,column sep=large]
 M_{0,3}^{\prime\sim}[0] \arrow[r] \& \Sigma_{\rho,\sigma}^{\text{\textcircled{eff}}}
\end{tikzcd}} \unsure{all}

b) $\forall u\in M_{0,3}^{\prime\sim},\ \exists u_{\unsure{\mathcal{G}}}\in\operatorname{Aut}\mathcal{G}$ such that
$u_{\unsure{M|\mathcal{G}}}\circ\psi=\psi\circ u$

b') Idem with $u\in M_{0,3}^{\prime\sim}[0]$

\unsure{c)} $\exists\psi_M$ which extends $\psi$

c') $\exists\psi_G$ which extends $\psi$

We will therefore say that \deleted{\unsure{effective}} \unsure{automatic|automorphism}

$\psi$ is \uline{\unsure{effective}} if (20) b) satisfied}

\uline{Definition} We will say that $\psi$ is \uline{effective} if for every
$u\in M_{0,3}^{\prime\sim}[0]$, putting $u(\rho)=\zeta\rho$, $u(\sigma)=\eta\sigma$,
\[
 (\zeta_{\mathcal{G}},\eta_{\mathcal{G}})
 =(u(\zeta),u(\eta))\in\Sigma_{\rho,\sigma}
\]
effective (i.e. $\in\Sigma_{\rho,\sigma}^{\mathrm{eff}}$, i.e. $\exists u_{\mathcal{G}}$ \unsure{automorphism} $\mathcal{G}$ such that we have \deleted{$\unsure{u_{\mathcal{G}}\rho}$}
\[
 (69)\qquad
 u_{\mathcal{G}}(\rho_{\mathcal{G}})=\zeta_{\mathcal{G}}\rho_{\mathcal{G}},
 \qquad
 u_{\mathcal{G}}(\sigma_{\mathcal{G}})=\eta_{\mathcal{G}}\sigma_{\mathcal{G}}
\]
or, what amounts to the same thing, that we have commutativity in
\diagram{\begin{tikzcd}[ampersand replacement=\&,column sep=large,row sep=large]
 \mathcal{G}_{0,3}^{+\wedge} \arrow[r,"u"] \arrow[d,"\psi"'] \& \mathcal{G}_{0,3}^{+\wedge} \arrow[d,"\psi"] \\
 \mathcal{G} \arrow[r,"u_{\mathcal{G}}"'] \& \mathcal{G}
\end{tikzcd}\quad(70).}

(Thus to say that this is the case for \unsure{every} $u\in M_{0,3}^{\prime\sim}\deleted{[0]}$ means that these $u$ transform $\operatorname{Ker}\psi$ into itself --- or also, that
$M_{0,3}^{\prime\sim}\deleted{\unsure{[0]}}$
$(=M_{0,3}^{\prime\sim}[0].\mathcal{G}_{0,3}^{+\wedge})$ in its entirety stabilizes $\operatorname{Ker}\psi$.
\deleted{When} This condition is therefore automatically fulfilled when $\mathcal{G},\rho_{\mathcal{G}},\sigma_{\mathcal{G}}$ satisfy \unsure{the condition|the conditions} b) of (20).

We have done what was needed to have \unsure{the}
\mspage{516}

\scanpage{518}
\mspage{562}

\uline{Theorem} Let $\psi=\psi_{\mathcal{G}}:\mathcal{G}_{0,3}^{+\wedge}\longrightarrow\mathcal{G}$ be an effective surjective homomorphism (cf. above).
Then there exists a unique homomorphism
\[
 (71)\qquad\psi_M:\deleted{\unsure{S_0}}M_{0,3}^{\prime\sim\deleted{\unsure{*}}}
 \longrightarrow\deleted{\unsure{S_0}}M_{\rho_{\mathcal{G}},\sigma_{\mathcal{G}}}^{0}
\]
\deleted{making the diagram commutative}
\diagram{\deleted{\begin{tikzcd}[ampersand replacement=\&,column sep=large,row sep=large]
 \mathcal{G}_{0,3}^{+\wedge} \arrow[r,"\unsure{\psi|\varphi}"] \arrow[d,"\psi"'] \& \deleted{\unsure{\mathcal{G}}}\,M_{0,3}^{\prime\sim} \arrow[d,"\psi"] \\
 \deleted{\unsure{\mathcal{G}}}\,\mathcal{G} \arrow[r] \& M_{\rho_{\mathcal{G}},\sigma_{\mathcal{G}}}^{0}\ (\overset{\mathrm{dfn}}{=}M_{\psi}^{0})
\end{tikzcd}\quad(72)}}

\margin{In fact, we have \unsure{of course in} the first place a homomorphism
\diagram{\begin{tikzcd}[ampersand replacement=\&,column sep=large]
 S_0M_{0,3}^{\prime\sim} \arrow[r] \& S_0M_{\psi}^{0}.
\end{tikzcd}}}

having the following property: if \deleted{$\unsure{u\in}$}
$u\unsure{=|\in}M_{0,3}^{\prime\sim}$ is of the form $u_{\alpha,\beta,f}$
($\alpha,\beta,f\in\mathcal{G}_{0,3}^{+\wedge}$, satisfying the loop equation
$\alpha\rho(\alpha)^{-1}=\beta\sigma(\beta)^{-1}\ell_1^{\nu}$), then
$\psi_M(u)$ is the element of \deleted{$\unsure{M_{\rho_{\mathcal{G}},\sigma_{\mathcal{G}}}^{0}}$} $M_{\psi}^{0}$ defined by the invariants $\psi(\alpha),\psi(\beta);\psi(f)$.

\uline{Corollary} We then have homomorphisms of exact sequences
\diagram{\begin{tikzcd}[ampersand replacement=\&,column sep=small,row sep=large]
 1 \arrow[r] \& \mathcal{G}_{0,3}^{+\wedge} \arrow[r,no head] \arrow[d,"\psi_{\mathcal{G}}"'] \& M_{0,3}^{\prime\sim} \arrow[r] \arrow[d,"\psi_M"] \& \Gamma_{\mathbb Q}^{\prime\sim} \arrow[r] \arrow[d,"\unsure{\psi_{\mathcal{T}}|\psi_{\Gamma}}"] \& 1 \\
 1 \arrow[r] \& \mathcal{G} \arrow[r] \& M_{\psi}^{0} \arrow[r] \& \deleted{\unsure{\mathcal{T}}}\Gamma_{\psi}^{0} \arrow[r] \& 1
\end{tikzcd}\quad(72).}

(The homomorphism $\psi_M$ is made explicit by means of three homomorphisms
\mspage{517}

\scanpage{519}

\[
\left\{\begin{aligned}
\psi_{T_\rho}&:\Gamma_{\mathbb Q}^{\prime\sim}\longrightarrow \unsure{Z_\rho|T_\rho}\\
\psi_{T_\sigma}&:\Gamma_{\mathbb Q}^{\sim}\longrightarrow \unsure{Z_\sigma|T_\sigma}
\end{aligned}\right.\qquad(73)
\]
\[
\psi_G:M_{0,3}^{\prime\sim}\longrightarrow G_{\rho,\sigma}\qquad(74)
\]
\margin{N.B. for simplicity we write $\rho,\sigma$ instead of $\rho_{\mathcal G},\sigma_{\mathcal G}$ as subscripts}

\deleted{satisfying the conditions}
(For \deleted{$\unsure{u}$} $x\in M_{0,3}^{\prime\sim}$,
\inserted{$\dot x$ its image in $\Gamma_{\mathbb Q}^{\prime\sim}$,}
\deleted{$a=\psi_{\unsure{T_\rho|\Gamma_\rho}}(\dot x),\ b=\psi_{\unsure{T_\rho|\Gamma_\rho}}(\unsure{y|\dot x}),\ U=\unsure{\psi u|\psi(x)}$}
\[
\deleted{\unsure{\delta_G(}}\quad
 a=\psi_{\unsure{T_\rho|\Gamma_\rho}}(\dot x),\quad b=\psi_{\unsure{T_\sigma|\Gamma_\sigma}}(\dot x),\quad U=\psi_G(x),
\qquad(75)
\]
we have
\[
\delta_\rho(a)=\delta_\sigma(b)=\delta_G(U)\quad\text{in }\gamma_{\rho,\sigma}\qquad(76)
\]
i.e. the three composite homomorphisms
\diagram{\begin{tikzcd}[ampersand replacement=\&,column sep=large,row sep=large]
\&\Gamma_{\mathbb Q}^{\prime\sim}\arrow[r,"\psi_{T_\rho}"]\&T_\rho\arrow[rd,"\delta_\rho"]\&\\
M_{0,3}^{\prime\sim}\arrow[ru,"\mathrm{can}"]\arrow[r,"\mathrm{can}"]\arrow[rrd,"\psi_G"',"\unsure{1}"{pos=.43}]\arrow[rrd,shorten >=55pt,shorten <=40pt]\&\Gamma_{\mathbb Q}^{\prime\sim}\arrow[r,no head,"\psi_{T_\sigma}"]\&T_\sigma\arrow[r,"\delta_\sigma"]\&\gamma_{\rho,\sigma}\arrow[ld,"\delta_G"]\\
\&\&G_{\rho,\sigma}\&
\end{tikzcd}\qquad(77)}
are equal. Moreover, we will have by (71),
\[
\psi_G|\mathcal G_{0,3}^{+\wedge}=\psi_{\mathcal G},\qquad(78)
\]
and finally this:

(79) If $x\in \deleted{\unsure{M}}M_{0,3}^{\prime\sim}$, $U=\psi_G(x)$, we have commutativity in
\diagram{\begin{tikzcd}[ampersand replacement=\&,column sep=huge,row sep=large]
\mathcal G_{0,3}^{+\wedge}\arrow[r,"\psi"]\arrow[d,"\operatorname{int}(x)|\mathcal G_{0,3}^{+\wedge}"']\&\mathcal G\arrow[d,"\operatorname{int}(U)|\mathcal G"]\\
\mathcal G_{0,3}^{+\wedge}\arrow[r,"\psi"]\&\mathcal G
\end{tikzcd}\qquad(80)}

\hfill $.\ .\ .\ .\ .\ .$ 
\mspage{518}

\scanpage{520}
\mspage{563}

. We can regard all this as a statement of \uline{functoriality} of the constructions made, with respect to a surjective homomorphism
\diagram{\begin{tikzcd}[ampersand replacement=\&]
\mathcal G\arrow[r,"\varphi=\varphi_{\mathcal G}"]\&\mathcal G'
\end{tikzcd}\qquad(81)}
\margin{\unsure{It would have} been better to take \unsure{then}
\diagram{\begin{tikzcd}[ampersand replacement=\&]\mathcal G'\arrow[r,"\varphi"]\&\mathcal G\end{tikzcd}},
in order to be able to specialize to $\mathcal G'$}
transforming the pair $(\rho,\sigma)$ of topological generators of $\mathcal G$ into a pair $(\rho',\sigma')$.
There is then a unique homomorphism
\diagram{\begin{tikzcd}[ampersand replacement=\&]
S_0M_{\mathcal G}^{0}\arrow[r,"\varphi_{M^0}"]\&S_0M_{\mathcal G'}^{0}
\end{tikzcd}\qquad(82)}
satisfying the following conditions:

\margin{this shows that $\varphi_{SM}$ is determined when we know the four homomorphisms $\varphi_B,\varphi_\rho,\varphi_\sigma,\varphi_G$}
a) $\unsure{\varphi_M|\varphi_{SM}}$ \unsure{passes} to quotients \unsure{as|to} homomorphisms
\[
\left\{\begin{aligned}
\varphi_B&:B_{\rho,\sigma}\longrightarrow\unsure{Z_{\rho,\sigma}|\Sigma_{\rho,\sigma}}\\
\varphi_\rho&:Z_\rho\longrightarrow Z_{\rho'}\\
\varphi_\sigma&:Z_\sigma\mathrel{\relbar\joinrel\relbar} Z_{\sigma'}\\
\varphi_G&:G_{\rho,\sigma}\mathrel{\relbar\joinrel\relbar} G_{\rho',\sigma'}
\end{aligned}\right.\ ;\qquad(83)
\]

b) \deleted{We have commutativity in the diagram,}
\inserted{\deleted{The homomorphisms $\varphi_B,\varphi_\rho,\varphi_\sigma$ \unsure{are induced by}}}
\[
(84)\quad\left\{\begin{gathered}
\deleted{\text{\unsure{only}}}\quad\varphi_G\qquad
\text{(which shows that knowledge}\\
\text{\illegible{1} of }\varphi_{SM}\ \unsure{b|se}\ \text{amounts to that of }\varphi_G\text{)};
\end{gathered}\right.
\]

c) We have commutativity in
\diagram{\begin{tikzcd}[ampersand replacement=\&,column sep=large,row sep=large]
\mathcal G\arrow[r]\arrow[d,"\varphi=\varphi_{\mathcal G}"']\&G_{\rho,\sigma}\arrow[d,"\varphi_G"]\\
\mathcal G'\arrow[r]\&G_{\rho',\sigma'}
\end{tikzcd}\qquad(85)}
(which shows, since $G_{\rho,\sigma}=B_{\rho,\sigma}\cdot\mathcal G$, that $\varphi_G$ is known when we know $\varphi_B$;

d) For $u\in B_{\rho,\sigma}$, \deleted{\unsure{denoting} by $u'$} \inserted{putting $u'=\varphi_B(u)$}, we have commutativity in
\diagram{\begin{tikzcd}[ampersand replacement=\&,column sep=large,row sep=large]
\mathcal G\arrow[r,"u"]\arrow[d,"\varphi_{\mathcal G}"']\&\mathcal G\arrow[d,"\varphi_{\mathcal G}"]\\
\mathcal G'\arrow[r,"u'"]\&\mathcal G'
\end{tikzcd}\qquad(86)}
\mspage{519}